% Part: lambda-calculus
% Chapter: introduction
% Section: lambda-computable

\documentclass[../../../include/open-logic-section]{subfiles}

\begin{document}

\olfileid{lam}{int}{cmp}
\olsection{\usetoken{S}{lambda definable} تابعې محاسبه کېدونکې دي}

\begin{thm}
\ollabel{thm:lambda-computable}
که يوه جزوي تابع~$f$ يو لامبډا ترم !!{lambda defined} کړې وي، نو
هغه محاسبه کېدونکې ده۔
\end{thm}

\begin{proof}
فرض کړئ تابع~$f$ يو لامبډا ترم~$X$ !!{lambda defined} کړې ده۔ د~$f$
د محاسبې يوه غير صوري کړنلاره بيان کړو۔ پر ننوت $m_0$،
\dots،~$m_{n-1}$، ترم $X \num m_0 \ldots \num
m_{n-1}$ وليکئ۔ يوه ونه جوړه کړئ: لومړے د اصلي ترم ټول يو-ګامي
راکمول وليکئ؛ تر هغو لاندې د هر يوه ټول يو-ګامي راکمول (يعنې د اصلي
ترم دوه-ګامي راکمول) وليکئ؛ او همداسې دوام ورکړئ۔ که هر کله يوې
عددنښې ته ورسېدئ، هماغه د ځواب په توګه ورکړئ؛ که ورنۀ رسېږئ، تابع
ناتعريفه ده۔

د چرچ اصل ته په رجوع سره ويلے شو چې دا تابع محاسبه کېدونکې ده۔ د
قضيې د ثبوت ښه لاره به دا وي چې دغه لټون‌کړنلاره په بازګشتي ډول
بيان کړو۔ مثلاً، د بنسټيزو بازګشتي تابعو او اړيکو يوه لړۍ تعريفولے
شو، لکه ``$\fn{IsASubterm}$''، ``$\fn{Substitute}$''،
``$\fn{ReducesToInOneStep}$''، ``$\fn{ReductionSequence}$''،
``$\fn{Numeral}$'' او داسې نور۔ بيا د~$f(m_0, \dots, m_{n-1})$
د محاسبې جزوي بازګشتي کړنلاره دا ده چې د يو-ګامي راکمول داسې لړۍ
ولټوو چې له $X \num{m_0} \dots \num{m_{n-1}}$ څخه پيل او پر يوه
عددنښه پاے ته رسېږي، او له هغې عددنښې سره سم عدد بېرته ورکړو۔
جزئيات اوږده او ستړي کوونکي دي، خو له دې پرته عادي دي۔
\end{proof}

\end{document}
